\SourceLeaf{287}
\hspace*{\parindent}Die Logarithmen der Dreiecksseiten in Toisen sind
\begin{center}
\begin{tabular}{lr}
Hohehagen-Inselsberg&\(4{,}6393865\)\\
Inselsberg-Brocken&\(4{,}7353929\)\\
Brocken-Hohehagen&\(4{,}5502669\)
\end{tabular}
\end{center}
\hspace*{\parindent}Die Azimute sind
\begin{center}
\begin{tabular}{lrrr}
&\multicolumn{3}{c}{Standpunkt Brocken}\\
Inselsberg&\(5^0\)&\(42'\)&\(22''\)\\
Hohehagen&58&49&8\\
&\multicolumn{3}{c}{Standpunkt Hohehagen}\\
Brocken&238&9&2\\
Inselsberg&324&23&1\\
&\multicolumn{3}{c}{Standpunkt Inselsberg}\\
Hohehagen&144&55&51\\
Brocken&185&35&21
\end{tabular}
\end{center}
\hspace*{\parindent}Dabei braucht man zwischen den Werten auf dem Sphäroid und denen auf der Kugel nicht zu unterscheiden. Bei den Logarithmen der Abstände zeigt sich erst in der achten oder neunten Dezimalstelle ein Unterschied, bei den Azimuten erst in den Tausendsteln einer Sekunde. Für unseren Zweck sind Logarithmen mit vier Dezimalstellen und Azimute in Minuten bereits genauer als nötig. Die Rechnung nach den obigen Formeln ergibt damit die folgenden Reduktionen. Sie sind mit ihren Vorzeichen zu den Azimuten auf dem Sphäroid zu addieren, um die Azimute auf der Kugel zu erhalten.
\begin{center}
\begin{tabular}{lr}
Brocken-Inselsberg&\(+0''00055\)\\
Brocken-Hohehagen&\(+0{,}00196\)\\
Hohehagen-Brocken&\(-0{,}00238\)\\
Hohehagen-Inselsberg&\(-0{,}00332\)\\
Inselsberg-Hohehagen&\(+0{,}00428\)\\
Inselsberg-Brocken&\(-0{,}00083\)
\end{tabular}
\end{center}
\Needspace{9\baselineskip}
\hspace*{\parindent}Damit ergeben sich für die Reduktion der Winkel des Dreiecks auf dem Sphäroid (zwischen den geodätischen Linien) auf die Winkel des Kugeldreiecks (zwischen Großkreisbögen) die folgenden Änderungen
